\documentclass[11pt]{article}

\usepackage{amsmath, amssymb, amsthm, geometry}
\usepackage{hyperref}
\usepackage{tikz-cd}
\usepackage{booktabs}
\geometry{margin=1in}

\title{Semiclassical Reconstruction of Riemann Surfaces from Bohrification\\
{\large (Paper II: Branch Sheaves and the Geometric Source of $R_*$)}}

\author{Zhou-Li Chen \\ co-nlang Research}
\date{May 2026}

\newtheorem{theorem}{Theorem}[section]
\newtheorem{proposition}[theorem]{Proposition}
\newtheorem{lemma}[theorem]{Lemma}
\newtheorem{corollary}[theorem]{Corollary}
\newtheorem{definition}[theorem]{Definition}
\newtheorem{remark}[theorem]{Remark}

\begin{document}

\maketitle

\begin{abstract}
Paper~I \cite{PaperI} introduces an affine geometric phase map $R_*$ that decomposes observed flat phase cohomology into a logical deviation and an objective \emph{geometric basepoint} $[\eta_{\mathrm{geo}}] \in \check{H}^1(M \setminus \Delta, U(1)_{\mathrm{flat}})$. This paper identifies the pure geometric source of this basepoint: the Riemann surface $\Sigma$ of the complexified classical action (formally, the étale space of the branch sheaf). The basepoint $[\eta_{\mathrm{geo}}]$ is the period character $\oint p\,dq/\hbar : H_1(M \setminus \Delta, \mathbb{Z}) \to U(1)$, an invariant of cohomological degree 1 defined independently of the set-level (degree 0) monodromy permutation representation. We verify this identification for $T^2$ and $S^2$, computing the monodromy and relating it to the Maslov phase on the 3-chart equatorial cylinder. The $\mathbb{R}^3$ case requires separate treatment due to the topological triviality of the base.
\end{abstract}

%------------------------------------------------------------------
\section{Introduction}
%------------------------------------------------------------------

Paper~I constructs a map
$R_*: H^1(\mathcal{C}^{\mathrm{BS}}(A), \underline{U(1)})
\to \check{H}^1(M, U(1))$
from the Bohrification topos to Čech cohomology, and verifies it in
$T^2$ and $S^2$. The present paper addresses: \emph{what is the precise geometric source of $[\eta_{\mathrm{geo}}]$?}

The answer is that $[\eta_{\mathrm{geo}}]$ arises directly from the Riemann surface of the multi-valued classical action. The Riemann surface of this multi-valued action (formally, the \'etale space of the branch sheaf) carries two parallel invariants: a degree-0 monodromy permutation representation and a degree-1 period character defined by integrating the action 1-form. The period character generates the \v{C}ech cocycle for $[\eta_{\mathrm{geo}}]$, independently of the sheet permutation structure. This provides a clean geometric interpretation of the basepoint without invoking the technicalities of presheaf colimits.

%------------------------------------------------------------------
\section{Mathematical Preliminaries}
%------------------------------------------------------------------

\subsection{Complexification and Branch Locus}

Let $(M, \omega)$ be a symplectic manifold. A \emph{complexification}
$M_\mathbb{C}$ is a complex manifold containing $M$ as a real submanifold,
such that holomorphic functions extend smooth functions on $M$.

For the examples:
\begin{itemize}
  \item $T^2_\mathbb{C} = (\mathbb{C}^*)^2$ (complexified torus).
  \item $S^2_\mathbb{C} = \mathbb{CP}^1$ (Riemann sphere).
  \item $\mathbb{R}^3_\mathbb{C} = \mathbb{C}^3$ (flat complex space).
\end{itemize}

For a multi-valued holomorphic function $\phi: M_\mathbb{C} \to \mathbb{C}$,
the \emph{branch locus} $\Delta \subset M_\mathbb{C}$ is the set of points
where $\phi$ is not single-valued (branch points, branch cuts, caustics).

\subsection{Branch Sheaf and Étale Space}

\begin{definition}[Branch sheaf]
\label{def:branch-sheaf}
The \emph{branch sheaf} $\mathcal{F}_\phi \in \mathbf{Sh}(M_\mathbb{C})$
assigns to each open $U \subset M_\mathbb{C}$ the set
\[
  \mathcal{F}_\phi(U) :=
  \bigl\{\,\phi_\alpha : U \to \mathbb{C}\text{ holomorphic}
    : \phi_\alpha \text{ is a single-valued branch of }\phi\text{ on }U
  \,\bigr\},
\]
with restriction maps given by analytic restriction.
\end{definition}

\begin{remark}[Étale space]
The \emph{étale space} $\mathrm{Et}(\mathcal{F}_\phi) \to M_\mathbb{C}$
is the disjoint union $\bigsqcup_{x \in M_\mathbb{C}} \mathcal{F}_\phi(x)$
of stalks, equipped with the topology generated by sections
$\phi_\alpha \in \mathcal{F}_\phi(U)$ as open sets
$\{(x, \phi_\alpha(x)) : x \in U\}$. For the branch sheaf,
$\mathrm{Et}(\mathcal{F}_\phi) \cong \Sigma$, the Riemann surface of $\phi$.

Étale spaces avoid the technicalities of presheaf colimits and
sheafification: the étale space is constructed directly from stalks
and sections.
\end{remark}

\subsection{Flat vs.\ Holomorphic Line Bundles}

A \emph{flat line bundle} over $M$ is classified by locally constant
$U(1)$-valued transition functions, i.e.\ a \v{C}ech 1-cocycle with
coefficients in the constant sheaf $U(1)$:
\[
  \check{H}^1(M, U(1)) \cong \mathrm{Hom}(\pi_1(M), U(1)).
\]

A \emph{holomorphic line bundle} over $M_\mathbb{C}$ has holomorphic
transition functions, classified by the Picard group
$\mathrm{Pic}(M_\mathbb{C})$.

%------------------------------------------------------------------
\section{BS-Localized Characters and Étale Construction}
\label{sec:setup}
%------------------------------------------------------------------

\subsection{BS-Localized Character Sets (Recap)}

As established in Paper~I, a Bohr--Sommerfeld character $\chi_j \in \mathrm{Char}_{\mathrm{BS}}(\mathcal{L}_j)$ corresponds to a branch $\phi_j$ of the complexified action via the semiclassical limit, yielding the $U(1)$-valued phase $e^{i\phi_j/\hbar}$. For a given quantum number, the set of such characters is discrete, with one character per branch of the action on $\mathcal{L}_j^\mathbb{C}$. We now organize these branches geometrically.

\subsection{Étale Space of the Branch Sheaf}

Rather than constructing a presheaf colimit and sheafifying, we work
directly with étale spaces.

\begin{definition}[Étale space of the branch sheaf]
\label{def:etale-char}
Define the étale space $\mathrm{Et}(\mathcal{F}_\phi) \to M_\mathbb{C}$ as follows:
\begin{enumerate}
  \item \textbf{Stalks:} For each $x \in M_\mathbb{C}$, the stalk
        $\mathcal{F}_\phi(x)$ is the set of germs of holomorphic branches
        $\phi_\alpha$ at $x$. Each germ is a local solution to the
        Hamilton--Jacobi equation. The stalk is a discrete set (the number
        of branches at $x$).

  \item \textbf{Topology:} A basis of opens is given by sections
        $\{(x, \text{germ of }\phi_\alpha \text{ at } x) : x \in U\}$
        for each branch $\phi_\alpha$ on $U \subset M_\mathbb{C}$.

  \item \textbf{Étale property:} The projection
        $\pi: \mathrm{Et}(\mathcal{F}_\phi) \to M_\mathbb{C}$ sending
        $(x, \text{germ}) \mapsto x$ is a local homeomorphism.
\end{enumerate}

\textbf{Note:} This is the standard étale space construction for a sheaf.
It is a \emph{covering space} (or branched cover, near branch points),
not a line bundle. Over $M_\mathbb{C} \setminus \Delta$, it is an
$n$-sheeted covering where $n$ is the number of branches.
\end{definition}

\begin{definition}[Holomorphic line bundle of phases]
\label{def:phase-bundle}
Define the holomorphic line bundle $L_\phi \to M_\mathbb{C}$ by:
\begin{enumerate}
  \item \textbf{Fibers:} $L_\phi(x) = \mathbb{C} \cdot e^{i\phi_\alpha(x)/\hbar}$, the complex line spanned by the phase factor $e^{i\phi_\alpha(x)/\hbar}$ for a branch $\phi_\alpha$ defined at $x$.

  \item \textbf{Transition functions:} On overlaps $U_\alpha \cap U_\beta$,
        the transition is
        \[
          g_{\alpha\beta} = \exp\!\Bigl(\tfrac{i}{\hbar}(\phi_\alpha - \phi_\beta)\Bigr)
          \in \mathcal{O}^*(U_\alpha \cap U_\beta),
        \]
        a holomorphic function (non-vanishing).
\end{enumerate}

\textbf{Important distinction:} On the full complexified base $M_\mathbb{C}$,
the transitions $g_{\alpha\beta}$ are \emph{holomorphic}. The bundle $L_\phi$ becomes flat (with locally constant transitions) only when restricted to the real slice $M \setminus \Delta \subset M_\mathbb{C}$. This is \emph{not} because holomorphic functions become constant, but precisely because the differences $\phi_\alpha - \phi_\beta$ are constrained to be constant along Lagrangian leaves by the Bohr--Sommerfeld condition (Lemma~3.1 of Paper~I).

\textbf{Note on \'Etale space vs.~Line Bundle:} The \'etale space $\mathrm{Et}(\mathcal{F}_\phi)$ is a covering space (discrete stalks); the line bundle $L_\phi$ is a $U(1)$-bundle (continuous fibers). These are different objects.
\end{definition}

\begin{remark}
The flat line bundle $L_\phi$ restricted to $M \setminus \Delta$ is precisely the
\emph{geometric basepoint} $[\eta_{\mathrm{geo}}]$ constructed in Paper~I. Its Čech cocycle $\{g_{\alpha\beta}\}$ defines a class
in $\check{H}^1(M \setminus \Delta, U(1)_{\mathrm{flat}})$. The étale space $\mathrm{Et}(\mathcal{F}_\phi)$
is the Riemann surface of the multi-valued action---a covering space whose
monodromy permutes the branches and generates this flat phase.
\end{remark}

%------------------------------------------------------------------
\section{Main Theorem: Riemann Surface and Period Character}
\label{sec:main}
%------------------------------------------------------------------

\begin{theorem}[Riemann surface identification over punctured base]
\label{thm:main}
Let $M \in \{T^2, S^2\}$ with Bohr--Sommerfeld foliation and complexified
action $\phi$. Let $\Delta \subset M_\mathbb{C}$ denote the branch locus
(caustics, poles). Over the punctured base $M_\mathbb{C} \setminus \Delta$,
the \'etale space $\mathrm{Et}(\mathcal{F}_\phi)$ is homeomorphic to the
Riemann surface $\Sigma$ restricted to this punctured base:
\[
  \mathrm{Et}(\mathcal{F}_\phi)|_{M_\mathbb{C}\setminus\Delta}
  \;\cong\;
  \Sigma|_{M_\mathbb{C}\setminus\Delta}.
\]
(The full Riemann surface $\Sigma$ is a branched cover including ramification over $\Delta$; extending the \'etale space to include these branch points requires completion of the stalks, which we leave separate from this regular-locus homeomorphism.)

\textbf{Key observations on the structure of $E_\phi := \mathrm{Et}(\mathcal{F}_\phi)$:}
\begin{enumerate}
\item \textbf{Discrete fibers:} For each $x \in M_\mathbb{C} \setminus \Delta$,
      the stalk $\mathcal{F}_\phi(x)$ consists of germs of holomorphic branches
      $\phi_\alpha$ at $x$. By Lemma~3.1 of Paper~I (quoted below), the
      branch differences $\phi_\alpha - \phi_\beta$ are locally constant
      modulo $2\pi\hbar$, so the $U(1)$-valued phases
      $e^{i\phi_\alpha/\hbar}$ form a \emph{discrete subset} of $U(1)$.

\begin{quote}
\textbf{Lemma 3.1 (Paper I).} Let $\phi_\alpha, \phi_\beta$ be two branches of the multi-valued classical action defined on overlapping open sets $U_\alpha \cap U_\beta \subset M \setminus \Delta$. Then the difference $\phi_\alpha - \phi_\beta$ is locally constant modulo $2\pi\hbar$. That is, for any connected component of the overlap, there exists an integer $n_{\alpha\beta} \in \mathbb{Z}$ such that
\[
  \phi_\alpha(x) - \phi_\beta(x) = 2\pi\hbar \cdot n_{\alpha\beta}
  \quad \text{for all } x \in U_\alpha \cap U_\beta.
\]
\end{quote}

\item \textbf{Covering space structure:} The projection
      $\pi: E_\phi \to M_\mathbb{C} \setminus \Delta$ is a local homeomorphism
      (\'etale property). Since the fibers are discrete, $E_\phi$ is an unbranched
      \emph{covering space} over $M_\mathbb{C} \setminus \Delta$.

\item \textbf{Sheet structure:} For $S^2$, the cover has two sheets
      (corresponding to branches $\phi_N$ and $\phi_S$), classified by
      the transition function $g_{NS} = -1 \in U(1)$ on the equatorial
      overlap. This matches the monodromy of the Riemann surface $\Sigma_l$
      for the rigid rotor action.
\end{enumerate}

\textbf{Explicit homeomorphism and naturality.}
The homeomorphism maps a point $(x, \text{germ}_x(\phi_\alpha)) \in E_\phi$ to $(x, \phi_\alpha(x)) \in \Sigma$.\linebreak As established in standard complex analysis (Forster~\cite[Ch.1]{Forster81}), the \'etale topology generated by local sections coincides with the analytic continuation topology of the Riemann surface.
This map is well-defined and bijective over $M_\mathbb{C} \setminus \Delta$ because each germ uniquely determines the branch value.\linebreak Furthermore, this map is \emph{natural} with respect to restrictions: restriction of a character $\chi_\alpha$ corresponds to analytic restriction of $\phi_\alpha$, which commutes with the projection maps to $M_\mathbb{C}$. Thus, it is a strict isomorphism of covering spaces over the punctured base.
\end{theorem}

\begin{theorem}[Period character from monodromy]
\label{thm:period-monodromy}
Let $\Sigma = \mathrm{Et}(\mathcal{F}_\phi)$ be the Riemann surface (for a fixed BS sector with quantum number $l$, we write $\Sigma_l$).
The monodromy representation of $\Sigma$ is:
\[
  \rho_{\mathrm{mon}}: \pi_1(M_\mathbb{C} \setminus \Delta)
  \longrightarrow \mathrm{Aut}(\mathcal{F}_\phi(x_0)),
\]
where $\mathcal{F}_\phi(x_0)$ is the stalk (set of germs) at a basepoint $x_0 \in M_\mathbb{C} \setminus \Delta$. This is the permutation group of the fiber of the covering space.

\textbf{Cohomological degree stratification.} The covering structure exhibits two parallel invariants stratified by cohomological degree:
\begin{itemize}
\item \textbf{Degree 0 (monodromy permutation):} The representation $\rho_{\mathrm{mon}}$
      acts on the \emph{set} of branches by permutations. For $S^2$ with
      2 sheets, $\mathrm{im}(\rho_{\mathrm{mon}}) \subseteq S_2 \cong \mathbb{Z}/2$.
      This captures the sheet-exchange structure but ignores geometric phase information.
\item \textbf{Degree 1 (period character):} The geometric period character
      $\rho_{\mathrm{per}}: \pi_1(M \setminus \Delta) \to U(1)$ is defined \emph{independently} of the monodromy permutation, via integration of the action 1-form $p\,dq/\hbar$.
      It factors directly through the abelianization $H_1(M \setminus \Delta, \mathbb{Z})$ because the 1-form $p\,dq$ is closed on the Lagrangian leaves ($d(p\,dq) = \omega|_{\mathcal{L}} = 0$) and its periods around contractible loops are integer multiples of $2\pi\hbar$ by the Bohr--Sommerfeld condition.
      This yields two parallel, independent invariants acting on different algebraic structures:
      \begin{itemize}
          \item[] $\rho_{\mathrm{mon}}: \pi_1(M \setminus \Delta) \longrightarrow \mathrm{Aut}(\mathcal{F}_\phi(x_0))$ \quad \textbf{(Degree 0: set-level)}
          \item[] $\rho_{\mathrm{per}}: H_1(M \setminus \Delta, \mathbb{Z}) \longrightarrow U(1)$ \quad \textbf{(Degree 1: cohomological)}
      \end{itemize}
      For $S^2$ in the maximal angular momentum sector ($m = \pm l$) where classical turning points coalesce at the poles such that $\Delta = \{\text{north pole, south pole}\}$, we have $H_1(\mathbb{CP}^1 \setminus \Delta) \cong \mathbb{Z}$ (generated by the equator loop), and the period map evaluates to $e^{i\pi(2l+1)} = -1$.
\end{itemize}

\textbf{Note:} The period character and the $S_n$ monodromy are parallel invariants of the Riemann surface. For 2-sheeted covers (like $S^2$), the torsion in $H_1$ coincidentally matches $S_2 \cong \mathbb{Z}/2$. However, for general covers ($n \geq 3$), $\rho_{\mathrm{per}}$ is defined independently of the non-abelian $S_n$ representation. This corrects the claim in early formulations of Paper~I that $R_*$ was the abelianization of $S_n$.
\end{theorem}

\begin{remark}[Trivial case: $T^2$]
For $T^2$, $\Delta = \emptyset$ and $\Sigma$ is a trivial cover
(disjoint union of sheets, one per quantum number, with no gluing).
The monodromy is trivial, so $\rho_{\mathrm{per}} = 1$ identically.
\end{remark}

\begin{proof}[Proof of Theorem~\ref{thm:period-monodromy}]
We verify for $M = S^2$ over $\mathbb{CP}^1 \setminus \Delta$.

\textbf{Action branches for $S^2$.} For the rigid rotor with quantum number $l$,
the action $\phi(\theta) = \int p\,dq$ has two distinct branches on $\mathbb{CP}^1$:
\begin{itemize}
  \item $\phi_N$: action computed from north pole reference.
  \item $\phi_S$: action computed from south pole reference.
\end{itemize}
The branches differ by the \emph{Maslov phase jump} at each pole:
\[
  \phi_N - \phi_S = \pi\hbar(2l+1)
  \quad\text{on the equator.}
\]
This phase jump corresponds to the Maslov index of the full polar cycle. As detailed in standard geometric quantization literature (e.g., Bates--Weinstein \cite[\S 7.4]{BW97} or de Gosson \cite[Ch.~8]{Gos11}), the Maslov index is determined by the winding number of the tangent plane in the Lagrangian Grassmannian. For the rigid rotor on $S^2$, the poles are caustic points, and a full polar cycle has Maslov index $\mu = 2$. The Bohr--Sommerfeld condition for this cycle is:
\[
  \oint p_\theta\, d\theta = 2\pi\hbar\left(l + \frac{\mu}{4}\right) = 2\pi\hbar\left(l + \frac{1}{2}\right) = \pi\hbar(2l+1).
\]
This entire action quantum is accumulated as a phase jump between the two branches across the equator.

\textbf{Riemann surface.} For a fixed $l$, the action $\phi$ defines a
Riemann surface $\Sigma_l$ as the branched cover of $\mathbb{CP}^1$ with:
\begin{itemize}
  \item Two sheets over $\mathbb{CP}^1 \setminus \{\mathrm{poles}\}$ (one for $\phi_N$, one for $\phi_S$).
  \item Branch points at the north and south poles (ramification index 2).
  \item Monodromy $\mathbb{Z}/2$ around each pole.
\end{itemize}

\textbf{Justification for two sheets:} The rigid rotor action
$\phi(\theta) = \int \sqrt{2IE - L_z^2/\sin^2\theta}\, d\theta$ involves
a square root with branch points at the classical turning angles
$\sin^2\theta_0 = m^2/(l(l+1))$. For any fixed quantum numbers $(l, m)$,
analytic continuation around these branch points exchanges the two
branches of the square root, yielding monodromy $\mathbb{Z}/2$ and a
2-sheeted cover. In the maximal angular momentum sectors ($m = \pm l$),
the turning points $\theta_0 = 0$ and $\pi - \theta_0 = \pi$ coalesce at
the poles. For the full rigid rotor action-angle monodromy, see
\cite{BW97, LS26}.

\textbf{\\'Etale space over punctured base and the 3-chart cover.} The étale space $E_\phi$ as
constructed (Definition~\ref{def:etale-char}) has:
\begin{itemize}
  \item Following Paper~I, the regular locus of $S^2$ retracts to the cylinder $S^1 \times \mathbb{R}$. We cover the equator with a 3-chart adapted good cover $\{U_1, U_2, U_3\}$. To make the action single-valued on the sheets, we introduce a longitudinal branch cut connecting the poles.
  \item The \'etale space $E_\phi$ over this punctured base is an unbranched 2-sheeted cover. The transition between the two sheets occurs exactly when traversing the branch cut.
\end{itemize}

\textbf{Homeomorphism over punctured base.} Away from the poles:
\[
  E_\phi|_{\mathbb{CP}^1 \setminus \Delta}
  \;\cong\;
  \Sigma_l|_{\mathbb{CP}^1 \setminus \Delta},
\]
both being unbranched 2-sheeted covers.
Placing the branch cut such that it transversely intersects exactly one overlap (e.g., $U_1 \cap U_2$)---a choice of cocycle representative---the transition function in the \'etale space evaluates to the Maslov phase jump $g_{12} = \exp(i\pi(2l+1)) = -1$, with $g_{23} = g_{31} = 1$.
Crucially, for a good cover of the equatorial cylinder, the triple intersection $U_1 \cap U_2 \cap U_3$ is empty. Thus, the \v{C}ech 1-cocycle condition $g_{12}g_{23}g_{31} = 1$ is vacuous.
However, to obtain a non-trivial element in $\check{H}^1$, we must verify that this 1-cochain is \emph{not} a coboundary. A 1-cochain $\{g_{\alpha\beta}\}$ is a coboundary if $g_{\alpha\beta} = f_\alpha f_\beta^{-1}$ for some local functions $f_\alpha: U_\alpha \to U(1)$. Since the cover is a good cover of $S^2 \setminus \Delta \simeq S^1 \times \mathbb{R}$, the sets $U_\alpha$ are contractible, so any local $U(1)$-valued function is homotopic to a constant. Thus coboundaries correspond to triples $(f_1, f_2, f_3) \in U(1)^3$ modulo global $f \in U(1)$, giving coboundaries of the form $(f_1f_2^{-1}, f_2f_3^{-1}, f_3f_1^{-1})$. The cochain $(-1, 1, 1)$ cannot be written in this form: if $-1 = f_1f_2^{-1}$ and $1 = f_2f_3^{-1} = f_3f_1^{-1}$, then $f_2 = f_3 = f_1$, but $-1 = f_1f_1^{-1} = 1$ is a contradiction. Hence $(-1, 1, 1)$ represents a non-trivial class in $\check{H}^1(S^2 \setminus \Delta, U(1)_{\mathrm{flat}}) \cong U(1)$.
The map $\chi_N \leftrightarrow \phi_N$, $\chi_S \leftrightarrow \phi_S$ establishes this homeomorphism explicitly, matching the gluing data of $E_\phi$ with the \v{C}ech cocycle of the geometric basepoint $[\eta_{\mathrm{geo}}]$.

\textbf{At the poles.} The Riemann surface $\Sigma_l$ has branch points
where the two sheets meet. The étale space $E_\phi$ as currently defined
does not include pole data. A complete construction would require adding
stalks at the poles and specifying their branching behavior.

\medskip
\textbf{Naturality.} Restriction to smaller opens commutes with the
identification on both sides: restriction of $\chi_\alpha$ corresponds
to analytic restriction of $\phi_\alpha$, both compatible with gluing.
The projection maps $E_\phi \to M_\mathbb{C}$ and $\Sigma \to M_\mathbb{C}$
intertwine with the homeomorphism, making it an isomorphism of covering
spaces.
\end{proof}

\begin{remark}[Consistency of phase distribution with Paper I]
In Paper~I \cite{PaperI}, the geometric phase was distributed symmetrically across the 3-chart cover ($\eta_{\mathrm{geo},\alpha\beta} = e^{i\pi/3}$ per overlap). In the current paper, we concentrate the phase jump at a single branch cut ($g_{12} = -1$, others $1$). Since the triple intersection $U_1 \cap U_2 \cap U_3$ is empty, there is no \v{C}ech cocycle condition to satisfy, and both choices represent the same cohomology class $[-1] \in \check{H}^1(M \setminus \Delta, U(1)_{\mathrm{flat}})$. The concentrated jump used here more clearly reflects the sheet transition geometry of the \'etale space.
\end{remark}

\section{Period Lattice Interpretation of $[\eta_{\mathrm{geo}}]$}
\label{sec:period}
%------------------------------------------------------------------

\begin{lemma}[Domain of period characters for $T^2$]
\label{lem:identification}
For $M = T^2$ with BS foliation, the full character group of the BS subposet (the domain of possible phases) is isomorphic to:
\[
  H^1(\mathcal{C}^{\mathrm{BS}}(T^2), \underline{U(1)})
  \;\cong\;
  \mathrm{Hom}(H_1(T^2, \mathbb{Z}), U(1))
  \;\cong\;
  (U(1))^2.
\]
However, the \emph{specific} geometric basepoint $[\eta_{\mathrm{geo}}]$ corresponds to the trivial element in this group, mapping every fundamental cycle to $\exp(\tfrac{i}{\hbar}\oint_{\gamma_k} p\,dq) = 1$.
\end{lemma}

\begin{proof}
For $T^2$, the singular locus is empty ($\Delta = \emptyset$), so the regular locus is the entire manifold $T^2$.
By Lemma~2.2 of Paper~I, the poset cohomology of the local observable algebras is isomorphic to the flat \v{C}ech cohomology of the base:
\[
  H^1(\mathcal{C}^{\mathrm{BS}}(T^2), \underline{U(1)}) \cong \check{H}^1(T^2, U(1)_{\mathrm{flat}}).
\]
By standard topology, flat line bundles on $T^2$ are classified by their holonomy representations, yielding the isomorphism:
\[
  \check{H}^1(T^2, U(1)_{\mathrm{flat}}) \cong \mathrm{Hom}(\pi_1(T^2), U(1)) \cong \mathrm{Hom}(H_1(T^2, \mathbb{Z}), U(1)) \cong (U(1))^2.
\]
The geometric basepoint $[\eta_{\mathrm{geo}}]$ is determined by the classical action phases. Because $T^2$ has no caustics ($\Delta = \emptyset$), the classical action is inherently single-valued (no branch points). The topological triviality of the bundle stems directly from $\Delta = \emptyset$, whereas the Bohr--Sommerfeld condition $\oint p\,dq = 2\pi\hbar n_k$ merely selects the specific quantized states within this trivial bundle. Consequently, the phase cocycle is identically $g_{\alpha\beta} = \exp(\frac{i}{\hbar}(\phi_\alpha - \phi_\beta)) = 1$. The period character maps every fundamental cycle to $1$, meaning $[\eta_{\mathrm{geo}}]$ is the trivial element in $(U(1))^2$.
\end{proof}

\begin{lemma}[Period character for $S^2$]
\label{lem:identification-S2}
For $M = S^2$ with BS foliation, the geometric basepoint $[\eta_{\mathrm{geo}}]$ corresponds to the equator holonomy in the flat cohomology:
\[
  [\eta_{\mathrm{geo}}] \;\in\; \check{H}^1(S^2 \setminus \Delta, U(1)_{\mathrm{flat}}) \cong U(1).
\]
The holonomy is always $-1 \in U(1)$ (independent of $l$ as a $U(1)$ element). This flat class extends via the clutching construction to a line bundle on the full $S^2$ with Chern class $c_1 = 2l+1 \in \mathbb{Z}$, which distinguishes different quantum sectors.

\textbf{Key distinction: holonomy vs.~Chern class.}
\begin{equation}
\text{Holonomy: } e^{i\pi(2l+1)} = -1 \;\;(\text{same for all } l \in \mathbb{Z})
\quad\text{vs.}\quad
\text{Winding number: } 2l+1 \;\;(\text{distinguishes sectors})
\end{equation}
Different winding numbers $2l+1$ map to the same $U(1)$ element $-1$ under $\exp: \mathbb{R} \to U(1)$, but yield distinct Chern classes $c_1 = 2l+1 \in \check{H}^2(S^2, \mathbb{Z})$ via the connecting homomorphism. The quantum number $l$ distinguishes prequantum bundles with the same equatorial holonomy but different topological degrees.

\textbf{Mechanism:} Since $S^2 \setminus \Delta \cong S^1 \times \mathbb{R}$, its first cohomology with $U(1)$ coefficients is $\mathrm{Hom}(\mathbb{Z}, U(1)) \cong U(1)$. The holonomy around the equator loop extracts the $U(1)$ phase directly.

The flat bundle on $S^2 \setminus \Delta$ extends to the full $S^2$ via the \emph{clutching construction}: view $S^2$ as two hemispheres $D_N \cup D_S$ glued along the equator $S^1$. The flat bundle over the cylinder $S^2 \setminus \Delta \simeq S^1 \times \mathbb{R}$ restricts to a flat bundle on a neighborhood of the equator with holonomy characterized by the clutching function $g_{NS} = e^{i\pi(2l+1)} = -1 \in U(1)$. The winding number of this clutching function, extracted from the phase exponent, is $2l+1$, yielding a line bundle $\mathcal{O}(2l+1)$ on the full sphere.

\textbf{Sign conventions.} Our constructions yield \emph{flat} line bundles (Paper~I's $g_{\alpha\beta}$ are locally constant), not holomorphic bundles. Following Paper~I, we define the Chern class via the connecting homomorphism $\delta: \check{H}^1(S^2, U(1)) \to \check{H}^2(S^2, \mathbb{Z})$ from the exponential sequence $0 \to \mathbb{Z} \to \mathbb{R} \to U(1) \to 0$. For the 3-chart equatorial cover, evaluating the \v{C}ech coboundary for the cocycle $g_{NS} = e^{i\pi(2l+1)}$ yields $\delta([g]) = 2l+1$. We refer to the standard geometric quantization literature (e.g., Bates--Weinstein~\cite[\S 7.4]{BW97}) for the full derivation showing that this flat equatorial holonomy correctly reconstructs the prequantum line bundle $\mathcal{O}(2l+1)$.

Equivalently, one may reduce the 3-chart cocycle $(-1, 1, 1)$ on the equatorial cylinder to the standard 2-chart clutching: the cocycle $g_{12} = -1$ on $U_1 \cap U_2$ restricts to the equatorial overlap as $g_{NS} = e^{i\pi(2l+1)}$. The Chern class $c_1 = 2l+1 \in \check{H}^2(S^2, \mathbb{Z})$ of the extended bundle is then recovered via the connecting homomorphism $\delta: \check{H}^1(S^2, U(1)) \to \check{H}^2(S^2, \mathbb{Z})$ of the exponential sequence on the full sphere.
\end{lemma}

\begin{corollary}
\label{cor:period}
Under the identification of Lemma~\ref{lem:identification} and~\ref{lem:identification-S2}, the geometric basepoint $[\eta_{\mathrm{geo}}]$ corresponds to the period/clutching character:

\begin{itemize}
  \item $T^2$: The basepoint $[\eta_{\mathrm{geo}}]$ is given by $[\gamma] \mapsto \exp(\tfrac{i}{\hbar}\oint_\gamma p\,dq) = 1$. It is the trivial character, consistent with the trivial bundle.
  \item $S^2$: The geometric basepoint $[\eta_{\mathrm{geo}}]$ is the element $[-1] \in U(1)$. For $l=0$, this holonomy corresponds to a prequantum bundle with Chern class $c_1 = 1$.
\end{itemize}
\end{corollary}

\begin{proof}
By Paper~I, the cocycle for $[\eta_{\mathrm{geo}}]$ is
$g_{\alpha\beta} = \exp(\tfrac{i}{\hbar}(\phi_\alpha - \phi_\beta))$.
The branch difference $\phi_\alpha - \phi_\beta$ integrates to
$\oint_\gamma p\,dq$ for a loop $\gamma$ connecting the branches.

The phase $\exp(\tfrac{i}{\hbar}\oint_\gamma p\,dq)$ is a well-defined
element of $U(1)$ by the exponential map.

For $T^2$, this gives the period character $H_1(T^2, \mathbb{Z}) \to U(1)$ via the exponential map $\mathbb{R}/2\pi\hbar\mathbb{Z} \cong U(1)$,
valued on $H_1 \cong \mathbb{Z}^2$.

For $S^2$, the cocycle $g_{NS}$ is classified by its winding number in
$\pi_1(U(1)) \cong \mathbb{Z}$, corresponding to the Chern class via
the exponential sequence $0 \to \mathbb{Z} \to \mathbb{R} \to U(1) \to 0$.
\end{proof}

\begin{remark}[Connection to Theorem~\ref{thm:main}]
The homeomorphism $E_\phi \cong \Sigma$ (over punctured base) identifies
the gluing data $g_{\alpha\beta} = \exp(\tfrac{i}{\hbar}(\phi_\alpha - \phi_\beta))$
with the period character. For a path $\gamma$ connecting branches:
\[
  \phi_\alpha - \phi_\beta = \oint_\gamma p\,dq,
\]
so the étale space gluing matches the period lattice.

The relationship can be summarized conceptually:
$R_*$ factors through the period character as follows:
\begin{itemize}
  \item Characters $s_j$ on BS algebras $\mathcal{A}_j$ determine phases $e^{i\phi_j/\hbar}$
        via the BS--Lagrangian leaf correspondence.
  \item The Riemann surface $\Sigma$ encodes the multi-valued action $\phi$
        as branches over $M_\mathbb{C} \setminus \Delta$.
  \item The \'etale space $E_\phi$ is homeomorphic to $\Sigma$ (Theorem~\ref{thm:main}).
  \item Periods $\oint_\gamma p\,dq/\hbar$ give the cohomology class
        in $\check{H}^1(M, U(1))$, factoring through $H_1(M \setminus \Delta, \mathbb{Z})$.
\end{itemize}
This gives the chain: BS subposet $\to$ Lagrangian leaves $\to$ Riemann surface
$\to$ \'etale space, with cohomological invariants tracking the period character.
\end{remark}

\begin{remark}[Correcting the $S_n$ monodromy claim]
Paper~I's original formulation incorrectly identified $R_*$ with
abelianization of a permutation representation $\rho: \pi_1 \to S_n$.
Theorem~\ref{thm:period-monodromy} above provides the correct interpretation
via cohomological degree stratification:
\begin{itemize}
\item Sheet permutation in $S_n$ captures the \emph{monodromy group}
      of the Riemann surface (how sheets are exchanged)---this is degree 0.
\item Phase cohomology $\check{H}^1(M, U(1))$ captures the
      \emph{abelian periods} $\oint p\,dq / \hbar$---this is degree 1.
\item These live in different cohomological degrees: monodromy in
      $H^0$ (covering structure), periods in $H^1$ (line bundle).
\item For 2-sheeted covers like $S^2$, the period character factors
      through $H_1(M \setminus \Delta, \mathbb{Z})$, which is the
      abelianization of $\pi_1$. The coincidence of $S_2 \cong \mathbb{Z}/2$
      with the torsion in $H_1$ is special to the 2-sheeted case.
\end{itemize}
For 2-sheeted covers like $S^2$, the permutation is $\mathbb{Z}/2$
and the period phase is $\pi(2l+1)$, giving $-1 \in U(1)$.
\end{remark}

%------------------------------------------------------------------
%------------------------------------------------------------------
\section{The $\mathbb{R}^3$ Case: Limitations and Sketch}
\label{sec:R3}
%------------------------------------------------------------------

The $\mathbb{R}^3$ example (harmonic oscillator) requires separate treatment.

\textbf{Topological triviality.} As a base space, $\mathbb{R}^3$ is
contractible, so $\check{H}^1(\mathbb{R}^3, U(1)) = 0$ and $R_* = 0$.
Every flat line bundle over $\mathbb{R}^3$ is trivial.

\textbf{Branch structure is NOT trivial.} However, the complexified action
for the harmonic oscillator has caustics (turning surfaces) and branch
points. The Riemann surface $\Sigma$ is a \emph{nontrivial branched cover}
of $\mathbb{C}^3$, with monodromy determined by loops around the turning
surfaces.

\begin{remark}
The claim that $\mathcal{F}_\phi = \{*\}$ (terminal sheaf) for $\mathbb{R}^3$
is incorrect. The correct statement:
\begin{itemize}
  \item $\check{H}^1(\mathbb{R}^3, U(1)) = 0$: trivial as a flat line bundle
        over the contractible base.
  \item The branch sheaf $\mathcal{F}_\phi$ is \emph{nontrivial} as a
        covering sheaf (it records the action's monodromy).
  \item The Kochen--Specker obstruction lives in $\ker(R_*)$ and is
        \emph{topology-independent}, confirmed by $\check{H}^1 = 0$.
\end{itemize}
For a complete analysis, one must specify a domain avoiding turning
surfaces, or acknowledge that the monodromy is captured by the branch
sheaf rather than Čech cohomology of the base.

\textbf{Sketch of correct analysis.} To obtain a meaningful cohomological analysis, one should:
\begin{enumerate}
  \item Remove the turning surface $\Delta \subset \mathbb{R}^3$ (caustics of the harmonic oscillator), yielding a non-contractible space $\mathbb{R}^3 \setminus \Delta$.
  \item Compute $H_1(\mathbb{R}^3 \setminus \Delta, \mathbb{Z})$, which is generated by loops linking the turning surfaces.
  \item Construct the period character $\rho_{\mathrm{per}}: H_1(\mathbb{R}^3 \setminus \Delta) \to U(1)$ via action integrals around these loops.
  \item Relate this to the monodromy of the Riemann surface over $\mathbb{C}^3$.
\end{enumerate}
The full analysis is left for future work.
\end{remark}

%------------------------------------------------------------------
\section{Summary and Limitations}
%------------------------------------------------------------------

The main results:

\begin{enumerate}
  \item \textbf{Theorem~\ref{thm:main}}: For $T^2$ and $S^2$, the étale space
        of BS-localized characters is homeomorphic to the Riemann surface
        of the classical action (for a fixed BS sector).
  \item \textbf{Corollary~\ref{cor:period}}: $[\eta_{\mathrm{geo}}]$ is the period/clutching
        character (via $H_1$ for $T^2$ and winding number for $S^2$),
        not abelianization of $S_n$ monodromy.
  \item The kernel as an open question: The nature of $\ker(R_*)$ for genuinely non-abelian monodromy remains to be determined.
\end{enumerate}

\textbf{Limitations:}
\begin{enumerate}
  \item \textbf{The kernel as an open question:} For compact integrable systems with
        genuinely non-abelian monodromy (degree $\geq 3$ covers with $S_n$ non-abelian),
        the exact nature of $\ker(R_*)$ remains an open question for future work.
        The examples $T^2$ and $S^2$ both have $\ker(R_*) = \{1\}$ (see table below);
        no concrete example with non-trivial kernel is presented here.
        Paper~III addresses the kernel via Kochen--Specker contextuality, but for a
        finite-dimensional algebra rather than classical monodromy.

  \item The $\mathbb{R}^3$ case requires specifying a domain avoiding
        turning surfaces; the branch sheaf is nontrivial.
  \item The homeomorphism in Theorem~\ref{thm:main} is for a fixed BS sector,
        not the full multi-valued action over all quantum numbers.
\end{enumerate}

\textbf{Summary table:} The verified and conjectural models:

\textbf{Terminology.} We use ``geometric basepoint'' to refer to the cohomology class $[\eta_{\mathrm{geo}}] \in \check{H}^1(M \setminus \Delta, U(1)_{\mathrm{flat}})$, and ``period character'' to refer to the homomorphism $\pi_1(M \setminus \Delta) \to U(1)$ representing it. These are equivalent data by the flat bundle classification $\check{H}^1(M, U(1)) \cong \mathrm{Hom}(\pi_1(M), U(1))$.

\begin{remark}[Triviality of $\ker(R_*)$ for $T^2$ and $S^2$]
For both $T^2$ and $S^2$, the affine geometric phase map $R_*$ is injective, so $\ker(R_*) = \{1\}$.
\begin{itemize}
\item For $T^2$: The domain of $R_*$ is $(U(1))^2$ (from Lemma~\ref{lem:identification}), and the geometric basepoint $[\eta_{\mathrm{geo}}]$ is the trivial element. The map $R_*$ is an isomorphism onto its image, so the kernel is trivial.
\item For $S^2$: The domain is $\check{H}^1(S^2 \setminus \Delta, U(1)) \cong U(1)$, a 1-dimensional Lie group (from Lemma~\ref{lem:identification-S2}). The map $R_*$ is affine with injective linear part (it maps the generator of $H_1$ to a non-trivial phase $e^{i\pi(2l+1)} = -1$). Since the domain is 1-dimensional and the image is non-trivial, $R_*$ is injective, so $\ker(R_*) = \{1\}$. The element $[\eta_{\mathrm{geo}}] = -1$ is the image of the generator, confirming the map is non-degenerate.
\end{itemize}
In both cases, the kernel is trivial because the geometric basepoint is uniquely determined by the classical action data.
\end{remark}

\begin{center}
\renewcommand{\arraystretch}{1.3}
\begin{tabular}{lccc}
\toprule
$M$ & Riemann surface structure & Chern class of $[\eta_{\mathrm{geo}}]$ & $\ker(R_*)$ (Logical) \\
\midrule
$T^2$ & trivial (1 sheet per $n$) & $c_1 = 0$ (trivial) & $\{1\}$ \\
$S^2$ & 2-sheeted, $\mathbb{Z}/2$ monodromy & $c_1 = 2l+1$ & $\{1\}$ \\
Higher $n \geq 3$ & $n$-sheeted, $S_n$ non-abelian & factors through $H_1$ & \textbf{Open Question} \\
\bottomrule
\end{tabular}
\end{center}

\textbf{Note:} We use multiplicative notation $\{1\} \subset U(1)$ for trivial
kernel. For $T^2$, the geometric basepoint is trivial; for $S^2$, the kernel
is trivial and the geometric basepoint extends to a prequantum bundle with Chern class $c_1 = 2l+1$.
This table covers compact examples only. For $\mathbb{R}^3$,
$\check{H}^1 = 0$ and the KS obstruction is topology-independent.

%------------------------------------------------------------------
\begin{thebibliography}{9}

\bibitem[BG81]{BG81}
Boutet de Monvel, L. and Guillemin, V.
\textit{The Spectral Theory of Toeplitz Operators}.
Ann.\ of Math.\ Studies~99, Princeton University Press, 1981.

\bibitem[Bry93]{Bry93}
Bryant, J.
Bundle gerbes and loop groups.
\textit{arXiv:hep-th/9308076}, 1993.

\bibitem[BW97]{BW97}
Bates, S. and Weinstein, A.
\textit{Lectures on the Geometry of Quantization}.
Berkeley Mathematics Lecture Notes, AMS, 1997.

\bibitem[For81]{Forster81}
Forster, O.
\textit{Lectures on Riemann Surfaces}.
Springer, 1981.

\bibitem[Gos11]{Gos11}
de Gosson, M.~A.
\textit{Symplectic Methods in Harmonic Analysis and in Mathematical Physics}.
Birkh{\"a}user, Basel, 2011.

\bibitem[GU88]{GU88}
Guillemin, V. and Uribe, A.
Circular symmetry and the trace formula.
\textit{Invent.\ Math.}, 96 (1989), 385--423.

\bibitem[HLS09]{HLS09}
Heunen, C., Landsman, N.~P., and Spitters, B.
Bohrification.
In \textit{Deep Beauty}, ed.\ H.~Halvorson,
Cambridge University Press, 2011, pp.~271--313.

\bibitem[Lan16]{Lan16}
Landsman, N.~P.
Bohrification: from classical concepts to commutative algebras.
\textit{arXiv:1601.02794}, 2016.

\bibitem[LS26]{LS26}
Lohmiller, W. and Slotine, J.-J.
On computing quantum waves exactly from classical action.
\textit{Proc.\ R.\ Soc.\ A}, 482.2336 (2026), 20250413.

\bibitem[Mas72]{Mas72}
Maslov, V.~P.
\textit{Th\'eorie des perturbations et m\'ethodes asymptotiques}.
Dunod, Paris, 1972.

\bibitem[Mur96]{Mur96}
Murray, M.~K.
Bundle gerbes.
\textit{J.\ London Math.\ Soc.}, 54 (1996), 403--416.

\bibitem[PaperI]{PaperI}
Chen, Z.-L.
From Contextuality to Phase Cohomology:
A Computable Bridge Between Bohrification and Geometric Quantization.
\textit{Zenodo, DOI: 10.5281/zenodo.20072818}, 2026.

\end{thebibliography}

\end{document}